\subsection{Complex measures}

DEFINITION 6. -- One calls complex measure on T every continuous linear form on the complex vector space \(\mathcal{K}_{\mathbf{C}}(\mathrm{T})\) . \(^{2}\)

The space \(\mathcal{M}_{\mathbf{C}}(\mathrm{T})\) of complex measures on T is thus the dual of the Hausdorff locally convex space \(\mathcal{K}_{\mathbf{C}}(\mathrm{T})\) .

If \(m\) is a complex measure on \(\mathrm{T}\) , its restriction to \(\mathcal{H}(\mathrm{T})\) is a vectorial measure on \(\mathrm{T}\) with values in \(\mathbf{C}\) (regarded as a vector space over \(\mathbf{R}\) );

m is determined by this restriction, since if \(f = f_{1} + if_{2} \in \mathcal{K}_{\mathbf{C}}(\mathrm{T})\) , the real part \(f_{1}\) and the imaginary part \(f_{2}\) of f are in \(\mathcal{K}(\mathrm{T})\) , and \(m(f) = m(f_{1}) + im(f_{2})\) . Conversely, for every vectorial measure \(m_{0}\) on T with values in C, the formula \(m(f) = m_{0}(f_{1}) + im_{0}(f_{2})\) defines a complex measure m, the only one on T whose restriction to \(\mathcal{K}(\mathrm{T})\) is \(m_{0}\) . We shall therefore henceforth identify a complex measure with its restriction to \(\mathcal{K}(\mathrm{T})\) ; such a measure m is of the form \(m = \mu_{1} + i\mu_{2}\) , where \(\mu_{1}\) and \(\mu_{2}\) are two real measures on T, which are called, respectively, the real part and the imaginary part of m. The support of m is the union of the supports of \(\mu_{1}\) and \(\mu_{2}\) . One knows that m is majorizable (No.~3, Cor. of Prop. 4); we shall call absolute value of m the positive measure \(|m|\) corresponding to the absolute value \(|x_{1} + ix_{2}| = \sqrt{x_{1}^{2} + x_{2}^{2}}\) on C. One has \(|m| = (\mu_{1}^{2} + \mu_{2}^{2})^{1/2}\) (No.~4, Remark following Prop. 9), \(^{3}\) and \(|\mu_{1}| \leqslant |m|\) , \(|\mu_{2}| \leqslant |m|\) , \(|m| \leqslant |\mu_{1}| + |\mu_{2}|\) ; moreover, m is a measure with base \(|m|\) , and one can write \(m = h \cdot |m|\) , where \(h \in \mathcal{L}_{\mathbf{C}}^{\infty}(|m|)\) and \(|h(t)| = 1\) locally almost everywhere for \(|m|\) (No.~4, Prop. 9). \(^{4}\) The supports of m and \(|m|\) are the same.

For every mapping \(\mathbf{f}\) of \(\mathrm{T}\) into a complex Banach space \(\mathbf{F}\) , essentially integrable with respect to \(|m|\) , one can define (No.~7) the integral of \(\mathbf{f}\) with respect to \(m\) (corresponding to the \(\mathbf{R}\) -bilinear mapping \((\mathbf{x},\lambda)\mapsto\lambda\mathbf{x}\) of \(\mathbf{F}\times\mathbf{C}\) into \(\mathbf{F}\) ), which will be denoted \(\int\mathbf{f}dm\) ; it follows at once from the uniqueness property of Prop. 11 that (with the preceding notations) \(\int\mathbf{f}dm=\int\mathbf{f}d\mu_{1}+i\int\mathbf{f}d\mu_{2}=\int\mathbf{f}hd|m|\) . We therefore have \(m(f)=\int fdm\) for \(f\in\mathcal{H}_{\mathbf{C}}(\mathrm{T})\) . We say that \(\mathbf{f}\) is essentially integrable with respect to \(m\) if it is so for \(|m|;^{5}\) mappings that are \(m\) -integrable, \(m\) -measurable, locally \(m\) -integrable, \(m\) -negligible or locally \(m\) -negligible are defined similarly. We write

\[
\mathcal {L} _ {\mathrm{F}} ^ {p} (\mathrm{T}, m) \left(\mathrm{resp.} \overline {{\mathcal {L}}} _ {\mathrm{F}} ^ {p} (\mathrm{T}, m), \mathrm{L} _ {\mathrm{F}} ^ {p} (T, m)\right)
\]

instead of \(\mathcal{L}_{\mathrm{F}}^{p}(\mathrm{T},|m|)\) (resp. \(\overline{\mathcal{L}}_{\mathrm{F}}^{p}(\mathrm{T},|m|)\) , \(\mathrm{L}_{\mathrm{F}}^{p}(\mathrm{T},|m|)\) ); these are complex vector spaces.

For f to be m-integrable (resp. essentially m-integrable), it is necessary and sufficient that f be integrable (resp. essentially integrable) with respect to each of the four measures \(\mu_{1}^{+}, \mu_{1}^{-}, \mu_{2}^{+}, \mu_{2}^{-}\) , by virtue of the inequalities between \(|m|, |\mu_{1}|, |\mu_{2}|\) and the relations \(|\mu_{k}| = \mu_{k}^{+} + \mu_{k}^{-}\) (Ch. V, §2, No.~2, Prop. 3 and its Cor. 1).

If f is essentially m-integrable (resp. m-integrable), then \(|f|\) is essentially \(|m|\) -integrable (resp. \(|m|\) -integrable), and it follows from Prop. 11 that

\[
\left| \int \mathbf {f} d m \right| \leqslant \int | \mathbf {f} | d | m |.\tag{7}
\]

Let F and G be two complex Banach spaces, u a continuous linear mapping of F into G. If f is an essentially m-integrable (resp. m-integrable) mapping of T into F, then \(u \circ f\) is essentially m-integrable (resp. m-integrable) and \(\int(u \circ f) dm = u (\int f dm)\) ; this follows at once from the foregoing and the analogous proposition for essentially \(|m|\) -integrable functions (Ch. IV, §4, No.~2, Th. 1 and Ch. V, §1, No.~3, Lemma and Def. 3).

Let \(m\) be a complex measure on \(T\) and let \(h\) be a locally \(m\) -integrable complex function. For every function \(f \in \mathcal{K}_{\mathbf{C}}(T)\) , the function \(fh\) is \(m\) -integrable and the mapping \(f \mapsto \int fh dm\) is a complex measure, denoted \(h \cdot m\) and called the measure with density \(h\) with respect to \(m\) . If \(m = g \cdot |m|\) , it is clear that \(h \cdot m = hg \cdot |m|\) ; since, moreover, \(|g(t)| = 1\) locally almost everywhere for \(|m|\) , for \(\mathbf{f}\) to be essentially integrable for \(n = h \cdot m\) it is necessary and sufficient that \(\mathbf{f}h\) be essentially \(m\) -integrable, in which case \(\int \mathbf{f} dn = \int (\mathbf{f}h) dm\) . Moreover, \(|h \cdot m| = |h| \cdot |m|\) . We again say that every measure of the form \(h \cdot m\) is a complex measure with base \(m\) ; two complex measures \(m, m'\) are said to be equivalent if each has a density with respect to the other, or, what amounts to the same, if \(m' = h \cdot m\) with \(h\) locally \(m\) -integrable and \(h(t) \neq 0\) locally almost everywhere for \(|m|\) . It is clear that \(m\) and \(|m|\) are equivalent and that, for \(m\) and \(m'\) to be equivalent, it is necessary and sufficient that \(|m|\) and \(|m'|\) be so.

If m and \(m'\) are two complex measures on T, and f is a function with values in a complex Banach space F, essentially integrable (resp. integrable) for both m and \(m'\) , then, for any complex numbers \(\lambda\) and \(\lambda'\) , f is essentially integrable (resp. integrable) for \(\lambda m + \lambda'm'\) , and

\[
\int \mathbf {f} d (\lambda m + \lambda^ {\prime} m ^ {\prime}) = \lambda \int \mathbf {f} d m + \lambda^ {\prime} \int \mathbf {f} d m ^ {\prime}.
\]

For, this follows from the Cor. of Prop. 11 of No.~7.

In addition, it follows from the definitions that

\[
\left| \lambda m + \lambda^ {\prime} m ^ {\prime} \right| \leqslant \left| \lambda \right| \cdot \left| m \right| + \left| \lambda^ {\prime} \right| \cdot \left| m ^ {\prime} \right|.
\]

One calls conjugate measure of a complex measure m the complex measure \(\overline{m}\) defined by \(\overline{m}(f) = \overline{m(\overline{f})}\) for \(f \in \mathcal{K}_{\mathbf{C}}(\mathrm{T})\) . If \(m = \mu_{1} + i\mu_{2}\) , where \(\mu_{1}\) and \(\mu_{2}\) are real measures, then \(\overline{m} = \mu_{1} - i\mu_{2}\) and \(|\overline{m}| = |m|\) ; if \(m = h \cdot |m|\) then \(\overline{m} = \overline{h} \cdot |m|\) . If f is an essentially m-integrable (resp. \(m\) -integrable) function with complex values, then \(\overline{f}\) is essentially \(\overline{m}\) -integrable (resp. \(\overline{m}\) -integrable) and

\[
\int \overline {{f}} d \overline {{m}} = \overline {{\int f d m}}.
\]

PROPOSITION 12. --- Let m be a complex measure on T, p and q conjugate exponents (Ch. IV, §6, No.~4). The bilinear form \((f,g)\mapsto\int fg dm\) is defined and continuous on the product \(\mathcal{L}_{\mathbf{C}}^{p}(m)\times\mathcal{L}_{\mathbf{C}}^{q}(m)\) ; the inequality \(\left|\int fg dm\right|\leqslant\mathrm{N}_{p}(f)\mathrm{N}_{q}(g)\) holds, and \(\mathrm{N}_{q}(g)\) is the norm of the continuous linear form on \(\mathrm{L}_{\mathbf{C}}^{p}(m)\) deduced from the linear form \(f\mapsto\int fg dm\) by passage to the quotient.

Moreover, if \(1 \leqslant p < +\infty\) , then every continuous linear form on the complex vector space \(\mathcal{L}_{\mathbf{C}}^{p}(m)\) is of the type \(f \mapsto \int fg dm\) , where \(g\) is a function in \(\mathcal{L}_{\mathbf{C}}^{q}(m)\) , whose class in \(\mathrm{L}_{\mathbf{C}}^{q}(m)\) is uniquely determined.

Since \(m = h \cdot |m|\) , where \(|h(t)| = 1\) locally almost everywhere, the first assertion follows at once from Hölder's inequality (Ch. IV, §6, No.~4, Th. 2); the second derives from Prop. 3 of Ch. IV, §6, No.~4. Finally, if \(u\) is a continuous linear form on \(\mathcal{L}_{\mathbf{C}}^{p}\) , its restriction to the (real) subspace \(\mathcal{L}^p\) of \(\mathcal{L}_{\mathbf{C}}^{p}\) is a continuous R-linear mapping of \(\mathcal{L}^p\) into \(\mathbf{C}\) ; if \(1 \leqslant p < +\infty\) , it is therefore of the type \(f \mapsto \int fg_1 d|m| + i \int fg_2 d|m|\) , where \(g_1\) and \(g_2\) belong to \(\mathcal{L}^q\) (Ch. V, §5, No.~8, Th. 4); whence the final assertion, on setting \(g = (g_1 + ig_2)h^{-1}\) .

